\documentclass[12pt]{article}
\usepackage[utf8]{inputenc}
\usepackage[T1]{fontenc}
\usepackage{hyperref}
\title{Talodi Language Resources for K-12 Curriculum Development}
\author{Compiled by Claude Sonnet 5 (claude-sonnet-5) for Andr\'as Kornai}
\date{}
\begin{document}
\maketitle

\section{Introduction}
Talodi (also called Jomang; ISO 639-3 tlo) is a Niger-Congo language of the Talodi group, spoken in the Nuba Mountains of South Kordofan, Sudan. UNESCO lists it as critically endangered. The figure of about 1,500 speakers dates from 1989, and many Talodi now speak Sudanese Arabic. Whether Talodi and Heiban form a single group is debated.

\section{Primary Educational Resources}
The searches found no primers, textbooks or other school materials in Talodi.

\section{Linguistic and Reference Materials}
\begin{itemize}
\item Wikipedia entry on the language: \url{https://en.wikipedia.org/wiki/Talodi_language}
\item Wikipedia entry on the Talodi--Heiban grouping: \url{https://en.wikipedia.org/wiki/Talodi%E2%80%93Heiban_languages}
\item Two Leiden University scholarly publications appeared in the results. I did not open them, so I cannot say whether they cover Talodi in detail: \url{https://scholarlypublications.universiteitleiden.nl/access/item%253A2716494/view} and \url{https://scholarlypublications.universiteitleiden.nl/access/item%3A2728068/view}
\end{itemize}

\section{Religious and Community Resources}
\begin{itemize}
\item Joshua Project people-group page for the Talodi in Sudan: \url{https://www.joshuaproject.net/people_groups/22479/SU}
\end{itemize}

\section{Digital and Audio Resources}
No audio, video or corpus resources were found.

\section{Curriculum Development Context}
Almost nothing is documented online, and the community is shifting to Arabic. Any curriculum work would begin with documentation and with the speakers themselves.
\end{document}